\begin{pastexam}[label={px:nsysu-112-dmalg-03}]{中山資工所資安組 112 學年度 第 3 題}
（計算）

{[10\%]} There are 65 boxes in an office and the size of each box is different. Assume that no two boxes
have the same size. How many boxes at least can be used together to pack a stuff that is smaller than the
smallest box (a box of larger size can contain a smaller one)?
\end{pastexam}

\begin{pastsoln}
答案（採用的讀法）：有 $2^{65} - 1$ 種用盒子包裝的方法。

\textbf{題意}。 原題的問句「How many boxes at least can be used together」字面上是「至少可以一起用幾個
盒子」，那答案只是 $1$（東西比最小的盒子還小，任何一個盒子都裝得下），不需要計算，不像是 $10$ 分的題目。
本題放在離散數學的計數題裡，合理的讀法是：「用\textbf{至少一個}盒子把東西包起來（大盒套小盒），一共有
幾種用法？」 下面照這個讀法作答。

\textbf{計數}。 東西比最小的盒子小，所以放得進任何一個盒子。 盒子大小兩兩不同，挑出任何一組盒子，都能
依大小排好、一個套一個，東西放在最小的那個裡面；而且套法只有一種（只能大的套小的）。 所以「一種用法」和
「一個非空的盒子子集合」一一對應。 $65$ 個盒子的子集合有 $2^{65}$ 個，扣掉空集合（題目要求至少用一個），
共
\[
2^{65} - 1 = 36\,893\,488\,147\,419\,103\,231
\]
種用法。 （每個盒子「用」或「不用」兩種選擇，$2^{65}$ 是乘法原理的結果。）
\end{pastsoln}
